\section{人工评估（Human Evaluation）}

\subsection{为什么人工评估是``金标准''}

人工评估被视为开放式任务的\textbf{金标准}（Gold Standard），原因有两个：
\begin{enumerate}
    \item 它直接反映了我们最终关心的东西——人类对文本质量的判断
    \item 它是开发新自动评估方法的\textbf{校准标准}——每个新的自动指标都需要与人工评估进行相关性验证
\end{enumerate}

在进行人工评估时，通常需要标注者从多个维度评价文本质量：流畅性（fluency）、连贯性（coherence）、常识性（common sense）、风格（style）、语法正确性（grammaticality）、冗余度（redundancy）等。

\subsection{人工评估的五大挑战}

\begin{figure}[H]
\centering
\includegraphics[width=0.95\textwidth]{slides-images/slide-028.jpg}
\caption{人工评估面临的多重挑战：标注者之间的不一致性、同一标注者的不一致性、不可复现性、只能评估精确率不能评估召回率、激励不对齐\protect\footnotemark}
\end{figure}
\footnotetext{Slides 第29页。}

\begin{warningbox}{人工评估的五大陷阱}
\begin{enumerate}
    \item \textbf{速度慢、成本高}：标注者需要仔细阅读和评判文本，远慢于自动指标
    \item \textbf{标注者间不一致}（Inter-annotator Disagreement）：即使 5 位研究者花 2--3 小时讨论并编写详细的标注规范，独立标注时也只有 \textbf{67\%} 的一致率（50\% 为随机水平）
    \item \textbf{同一标注者的不一致}（Intra-annotator Disagreement）：同一个人在不同时间对同一样本可能给出不同评价
    \item \textbf{不可复现}：一项分析 128 篇论文的研究发现，仅 \textbf{5\%} 的人工评估是可复现的——大多数论文没有提供足够的实验设计细节
    \item \textbf{激励不对齐}：众包工人的目标是最大化时薪，而非最大化标注质量。AlpacaFarm 项目中发现工人速度是研究者的 2--3 倍，他们倾向于使用捷径（如偏好更长的回答）
\end{enumerate}
\end{warningbox}

\begin{knowledgebox}{人工评估只能评估``精确率''不能评估``召回率''}
给定模型的一个生成结果，人类可以评估\textbf{这个}结果的质量（精确率视角）。但人类无法评估模型\textbf{所有可能生成}的质量——那需要大量采样，时间和成本都不可承受。这是人工评估的一个根本性局限。
\end{knowledgebox}

\subsection{实践中的额外挑战}

进行人工评估时需要考虑的实操问题：

\begin{itemize}
    \item \textbf{任务描述}：必须提供非常详细的标注规范（rubric）
    \item \textbf{展示方式}：两个回答的左右/上下位置会影响标注者的选择
    \item \textbf{标注者筛选}：需要先用已知答案的样本测试标注者的能力
    \item \textbf{持续监控}：在每批标注中插入``金标准''样本，检测标注质量是否退化；监控标注速度（速度突然加快可能意味着偷懒）
    \item \textbf{绝不能跨论文比较}：不同论文使用不同的标注者、不同的标注规范、不同的展示方式，得到的人工评估结果完全不可比
\end{itemize}

\subsection{Chatbot Arena：大规模人工评估}

\begin{figure}[H]
\centering
\includegraphics[width=0.95\textwidth]{slides-images/slide-032.jpg}
\caption{Chatbot Arena 工作原理：用户在线免费使用模型，两个模型并排展示对同一问题的回答，用户选择偏好的一方，最终通过 Elo 评分系统生成排行榜\protect\footnotemark}
\end{figure}
\footnotetext{Slides 第33页。}

Chatbot Arena 是当前评估聊天 LLM 的\textbf{最主流人工评估基准}。其核心思路是借鉴国际象棋的 Elo 评分系统：

\begin{itemize}
    \item 任何人都可以上线免费使用各种顶尖模型
    \item 用户只需选择``左边更好''或``右边更好''
    \item 积累了约 \textbf{200,000} 条人类投票后，通过锦标赛式的 Elo 评分生成排行榜
    \item 不需要每对模型都直接对比——Elo 系统可以从部分比较中推断全局排名
\end{itemize}

\begin{warningbox}{Chatbot Arena 的局限性}
\begin{itemize}
    \item 用户是自选的``网上随机用户''，其偏好可能不具代表性（但数据量够大时这个问题会缓解）
    \item \textbf{成本极高}：需要巨大的社区努力和大量用户参与
    \item 只有知名模型（OpenAI、Anthropic、Google、Meta 等）才能获得足够多的评估——你自己训练的小模型永远不会有 200,000 人免费为你标注
    \item 即使是大公司，也无法将其用于日常开发中的超参调优——它只能用于最终的模型选择
\end{itemize}
\end{warningbox}

\subsection{本章小结}

人工评估是开放式任务的金标准，但面临速度慢、成本高、标注者不一致、不可复现、激励不对齐等严峻挑战。Chatbot Arena 通过大规模众包和 Elo 评分系统部分解决了这些问题，但其高成本和只适用于知名模型的局限性仍然存在。这些问题推动了 LLM 作为评估者（LLM-as-Judge）的研究方向。

%% ========== Part 5: 无参考评估与 LLM 作为评估者 ==========

